\documentclass[11pt]{article}

% --- Series style ---
\usepackage{series}

% --- Math and layout ---
\usepackage{amsmath,amssymb,amsthm}
\usepackage{mathtools}
\usepackage{geometry}
\usepackage{hyperref}

\geometry{margin=1in}


% --- Metadata ---
\title{Geometry--Light Relations in Modal Triplet Theory:\\Exact Identities, Conditional Bounds, and Principal Symbols}
\author{Peter Nero}
\date{July 2026}

\begin{document}
\maketitle

\begin{abstract}
We separate the geometry-light statements of MTT into exact algebraic
identities, consequences of explicit symmetry or positivity assumptions, and
phenomenological estimates requiring additional dynamics. Modal democracy
implies $\sin^2\theta_W=3/8$, but democracy is an assumption and is not selected
by the current gauge profile. Holonomy phase sums require a specified trivial
tensor product and compatible connection. Equal propagation speeds require
equality of principal symbols and do not follow from topology alone.
Curvature--mass drift is an identity conditional on a chosen mass--curvature
ansatz, while the PPN statement requires an explicit response-norm estimate.
No internal spectral gap is identified with an external cutoff. These typed
statements provide safe interfaces for later numerical tiers without
overstating their physical content.
\end{abstract}

\section{Classification rule}

A relation is called \emph{exact} only when it follows algebraically from its
displayed definitions. A \emph{conditional theorem} additionally lists every
symmetry, bundle, principal-symbol, positivity, or response assumption. A
\emph{phenomenological bound} also requires a selected physical normalization
and comparison scheme. Geometry-light does not mean assumption-free.

\section{Gauge-weight identity and modal democracy}

At one common scale and in one scheme, write
\[
\alpha_r^{-1}=K\zeta_r,
\qquad \zeta_r>0,
\]
and use $g_Y^2=(3/5)g_1^2$. Then the exact algebraic identity is
\begin{equation}
\boxed{\sin^2\theta_W
=\frac{1}{1+(5/3)(\zeta_1/\zeta_2)}.}
\label{eq:weak-general}
\end{equation}

\begin{definition}[Modal democracy]
Modal democracy at the chosen scale means $\zeta_1=\zeta_2$.
\end{definition}

\begin{corollary}[Conditional democracy value]
If modal democracy holds, Equation~\eqref{eq:weak-general} gives
$\sin^2\theta_W=3/8$.
\end{corollary}

This is exact given democracy, but democracy is not derived by this paper and
is not satisfied by the selected low-scale SMDR profile. The value $3/8$ is
therefore neither a current MTT weak-angle prediction nor evidence for a
physical matching scale. For
$\zeta_1/\zeta_2=1+\varepsilon$, the Taylor identity is
\[
\sin^2\theta_W=\frac38\left(1-\frac58\varepsilon
+\frac{25}{64}\varepsilon^2+O(\varepsilon^3)\right).
\]

\section{Conditional holonomy phase sum}

\begin{theorem}[Holonomy product under compatible trivialization]
Let $L_{12},L_{23},L_{31}$ be line bundles equipped with unitary connections.
Assume a specified connection-preserving trivialization
\[
L_{12}\otimes L_{23}\otimes L_{31}\cong\underline{\mathbb C}.
\]
Then for every closed loop $\gamma$,
\[
U_{12}(\gamma)U_{23}(\gamma)U_{31}(\gamma)=1,
\]
and the three phases sum to an integer multiple of $2\pi$.
\end{theorem}

\begin{proof}
The tensor-product connection is trivial by assumption, so its holonomy is
unity. Holonomy on a tensor product is the product of the three line-bundle
holonomies.
\end{proof}

Topological triviality alone does not force an arbitrary chosen product
connection to have zero holonomy. The compatible trivialization is an essential
hypothesis. The theorem constrains phase sums but does not select CKM or PMNS
phases individually.

\section{Principal-symbol condition for propagation speeds}

For a linearized field component $\phi_a$, let the second-order operator have
principal symbol
\[
\sigma_2(P_a)(x,\xi)=G_a^{\mu\nu}(x)\xi_\mu\xi_\nu.
\]

\begin{theorem}[Common characteristic cone]
Two canonically normalized sectors have the same local high-frequency
propagation cone if and only if their nondegenerate principal symbols define
the same null cone. In dimension at least three this means
$G_b^{\mu\nu}=\Omega^2G_a^{\mu\nu}$ for a positive conformal factor, subject to
the usual regularity assumptions.
\end{theorem}

Lower-order masses, potentials, projectors, or internal overlap coefficients do
not determine the characteristic cone. Equal wave speeds therefore require a
common principal-symbol theorem; they cannot be inferred solely from bundle
topology, modal counting, or an internal spectral gap.

\section{Representation-only RG signs}

Conditional on importing the Standard Model field content and perturbative
quantization, the one-loop coefficient
\[
b_0=\frac{11}{3}C_A-\frac{4}{3}\sum_fT(R_f)
-\frac{1}{6}\sum_sT(R_s)
\]
fixes the qualitative signs of the gauge beta functions. This is a standard
QFT consequence of representation content, not a new geometric prediction of
MTT. Quantitative running uses the selected multi-loop transport rather than
this sign-level approximation.

\section{Conditional curvature--mass drift}

Assume, rather than derive, the local ansatz
\[
m(x)=\sqrt{\kappa(\lambda^{(0)}+\beta R(x))},
\qquad \lambda^{(0)}+\beta R(x)>0.
\]
Then differentiation gives the exact identity
\[
\nabla_\mu\log m
=\frac{\beta\nabla_\mu R}{2(\lambda^{(0)}+\beta R)}.
\]
If the denominator is bounded below by $\lambda_{\min}>0$, then
\[
|\nabla_\mu\log m|
\leq\frac{|\beta|}{2\lambda_{\min}}|\nabla_\mu R|.
\]
The calculus is exact, but the mass ansatz, $\beta$, and physical-unit bridge
are independent inputs. This is not a prediction of cosmological mass drift.

\section{Conditional PPN estimate}

Suppose a linearized gravitational equation can be written
\[
\square\bar h_{\mu\nu}=-16\pi G_NT_{\mu\nu}+S_{\mu\nu}
\]
and a specified solution norm obeys
\[
\|S\|\leq C\Delta_{\mathrm{curv}}\|T\|.
\]
Only under a stable inverse estimate for the gauge-fixed operator does this
imply a bound of the form
\[
|\gamma-1|\leq C_{\mathrm{PPN}}\Delta_{\mathrm{curv}}.
\]
The constants and norm must be computed for the selected background and source.
Writing merely $\gamma=1+O(\Delta_{\mathrm{curv}})$ is an asymptotic template,
not a numerical solar-system prediction.

\section{No internal-gap cutoff inference}

An internal Laplacian gap controls suppression of omitted internal modes in a
specified reduction. It does not by itself define a four-dimensional UV,
coherence, inflationary, or cosmological cutoff. Any such identification needs
a dimensionful bridge, an action-level decoupling theorem, and an error bound.
The common SMDR point $Q=M_t$ is likewise a matching convention rather than a
physical cutoff.

\section{Conclusion}

The geometry-light tier contains useful exact relations, but their scope is
now explicit. The weak-angle and curvature formulas are algebraic identities
conditional on stated ansatz data; holonomy phases require a compatible
trivialization; wave-speed equality is a principal-symbol question; RG signs
are imported Standard Model QFT; and PPN control needs a quantitative response
estimate. These statements can constrain later constructions without being
misreported as source-derived numerical predictions.

\begin{thebibliography}{99}

\bibitem{MTT-Admissibility}
P.~Nero,
\emph{Modal Triplet Theory: Admissibility, Encodings, and the Structure of Physical Description},
Zenodo preprint, January 2026.
\url{https://doi.org/10.5281/zenodo.18255621}

\bibitem{MTT-Foundation}
P.~Nero,
\emph{Modal Triplet Theory: Foundation},
Zenodo preprint, September 2025.
\url{https://doi.org/10.5281/zenodo.16949762}

\bibitem{MTT-FixedPoints}
P.~Nero,
\emph{Fixed Points I--VI: Complete Coherence Spine},
Zenodo preprints, August 2025.
\url{https://doi.org/10.5281/zenodo.16948748}

\bibitem{MTT-ProjectionAdmissibility}
P.~Nero,
\emph{The Projection--Admissibility Principle: Structural Constraints on Effective Physical Description},
Zenodo preprint, January 2026.
\url{https://doi.org/10.5281/zenodo.18255838}

\bibitem{MTT-Closure}
P.~Nero,
\emph{Closure and Inevitability in Modal Triplet Theory},
Zenodo preprint, January 2026.
\url{https://doi.org/10.5281/zenodo.18255510}

\bibitem{MTT-CoherenceCapacity}
P.~Nero,
\emph{Coherence Capacity as the Fundamental Resource of Effective Physics},
Zenodo preprint, January 2026.
\url{https://doi.org/10.5281/zenodo.18255905}

\bibitem{MTT-CoherenceDynamics}
P.~Nero,
\emph{Dynamics of Coherence Capacity: Transport, Concentration, and Exhaustion},
Zenodo preprint, January 2026.
\url{https://doi.org/10.5281/zenodo.18256048}

\bibitem{MTT-QM}
P.~Nero,
\emph{Modal Triplet Theory: From MTT to Quantum Mechanics},
Zenodo preprint, September 2025.
\url{https://doi.org/10.5281/zenodo.17074246}

\bibitem{MTT-QFT}
P.~Nero,
\emph{From MTT to Quantum Field Theory},
Zenodo preprint, 2025.
\url{https://doi.org/10.5281/zenodo.17068816}

\bibitem{MTT-GR}
P.~Nero,
\emph{Modal Triplet Theory: From MTT to General Relativity},
Zenodo preprint, October 2025.
\url{https://doi.org/10.5281/zenodo.16950597}

\bibitem{MTT-UVQG}
P.~Nero,
\emph{Modal Triplet Theory: From MTT to a UV-Finite, Unitary Quantum Gravity},
Zenodo preprint, 2025.
\url{https://doi.org/10.5281/zenodo.17077671}

\bibitem{MTT-Measurement}
P.~Nero,
\emph{Measurement as Disturbance and Stabilization in Modal Triplet Theory},
Zenodo preprint, 2025.
\url{https://doi.org/10.5281/zenodo.17177404}

\bibitem{MTT-Irreversibility}
P.~Nero,
\emph{Projection, Probability, and Irreversibility: Shadow Bridges Between Measurement, Black Holes, and Cosmology in Modal Triplet Theory},
Zenodo preprint, January 2026.
\url{https://doi.org/10.5281/zenodo.18256408}

\bibitem{MTT-Bell-Beables}
P.~Nero,
\emph{Modal Fixed Points, Bell’s Beables, and the Limits of Factorization},
Zenodo preprint, 2025.
\url{https://doi.org/10.5281/zenodo.17076300}

\bibitem{MTT-Bell-Temporal}
P.~Nero,
\emph{Temporal Bell Inequalities and Global Consistency in Modal Triplet Theory},
Zenodo preprint, August 2025.
\url{https://doi.org/10.5281/zenodo.18208884}

\bibitem{MTT-Stochastic}
P.~Nero,
\emph{From Modal Triplet Theory to Indivisible Stochastic Processes: A First-Principles, Fully Rigorous Derivation},
Zenodo preprint, January 2026.
\url{https://doi.org/10.5281/zenodo.18254862}

\end{thebibliography}
\end{document}
